\startAnhang

\listofanhang
\clearpage

% Jeder Anhang wird mit \anhang{Bezeichnung} begonnen.
% Unterteilungen mit \anhangteil{Bezeichnung}.
%
% Falls Sie (noch) keinen Anhang benötigen, können Sie die Einbindung
% von includes/anhang.tex in der Hauptdatei auskommentieren.

\anhang{Gesprächsprotokoll zum Nutzungskontext}
\label{anhang:gespraech}

Gespräch mit Benjamin Kaiser, Meister Instandhaltung, Woodward L'Orange
GmbH, am 10.06.2026 in Glatten, Dauer: ca. 30 Minuten. Gegenstand:
Benutzergruppe, Nutzungsfrequenz und Umgebungsbedingungen der Anwendung
\enquote{Meldung -- Foto aufnehmen}. Das nachfolgende Stichpunktprotokoll
wurde im Anschluss an das Gespräch durch den Verfasser erstellt.

\bigskip
\textbf{Stichpunktprotokoll:}
\begin{itemize}
  \item Nutzergruppe: ca. 20 Instandhaltungstechniker (geringe SAP-Erfahrung)
  \item SAP-spezifische Begriffe und Schlüssel in den Meldungen sind der
        Nutzergruppe in der Regel nicht geläufig
  \item Einführung der Anwendung auf Wunsch des Meisters Instandhaltung
  \item Nutzung freiwillig, liegt im Ermessen der einzelnen Mitarbeitenden
  \item Nutzungsfrequenz: gruppenweit täglich, individuell ca.\ zweimal
        wöchentlich
  \item Umgebung: Lärm, Schmutz, Ölrückstände, beengte Platzverhältnisse;
        Schutzbrille durchgängig, Handschuhe zeitweise
  \item Für die Dokumentation haben die Mitarbeitenden genug Zeit
        
\end{itemize}


\anhang{Kriterienkatalog zur heuristischen Evaluation}
\label{anhang:kriterienkatalog}

Der nachfolgende Katalog überführt die anwendbaren Empfehlungen der
ISO~9241-110:2020 (Tabelle~A.1) in Prüfkriterien. Alle Abschnittsangaben
beziehen sich auf diese Norm. Die Spalte \enquote{Befund} dokumentiert den
Inspektionsbefund (erfüllt/Abweichung) und den \ac{SG} nach
Tabelle~\ref{tab:schweregrade}.

{\footnotesize
\begin{longtable}{|>{\centering\arraybackslash}m{1cm}|>{\centering\arraybackslash}m{7cm}|>{\centering\arraybackslash}m{1.6cm}|>{\centering\arraybackslash}m{3.3cm}|>{\centering\arraybackslash}m{0.8cm}|}
\caption{Kriterienkatalog zur heuristischen Evaluation der Anwendung}
\label{tab:katalog} \\\hline
\textbf{Nr.} & \textbf{Prüfkriterium} & \textbf{ISO-Abs.} & \textbf{Befund} & \textbf{SG} \\\hline
\endfirsthead
\multicolumn{5}{c}{\tablename~\thetable{} -- Fortsetzung} \\\hline
\textbf{Nr.} & \textbf{Prüfkriterium} & \textbf{ISO-Abs.} & \textbf{Befund} & \textbf{SG} \\\hline
\endhead
\hline \multicolumn{5}{r}{\footnotesize Fortsetzung nächste Seite} \\
\endfoot
\hline
\endlastfoot

\multicolumn{5}{|l|}{\textbf{Prinzip 1 -- Aufgabenangemessenheit}} \\\hline
1.1 & Macht die Anwendung hinreichend deutlich, dass sie der Fotodokumentation zu einer bereits bestehenden Meldung dient, sodass Benutzende ihre Eignung erkennen? & 5.1.2 & erfüllt & 0 \\\hline
1.2 & Stehen in jedem Schritt des Ablaufs (Meldungsauswahl, Bestätigung, Fotoaufnahme/-upload, Fotoverwaltung) die erforderlichen Steuerelemente und Informationen bereit? & 5.1.3.1 & erfüllt & 0 \\\hline
1.3 & Verzichtet die Anwendung auf Funktionen und Informationen, welche die Fotodokumentation behindern oder ablenken? & 5.1.3.3 & erfüllt & 0 \\\hline
1.4 & Bietet die Anwendung dort, wo zweckmäßig, sinnvolle Vorbelegungen an (z.\,B. bei den Filterkriterien der Meldungssuche)? & 5.1.4.1 & erfüllt & 0 \\\hline
1.5 & Vermeidet die Anwendung Vorbelegungen, die zur Auswahl einer falschen Meldung oder zur Fehlzuordnung der Fotos führen können? & 5.1.4.2 & erfüllt & 0 \\\hline

\multicolumn{5}{|l|}{\textbf{Prinzip 2 -- Selbstbeschreibungsfähigkeit}} \\\hline
2.1 & Ist erkennbar, an welcher Stelle des Ablaufs sich Benutzende befinden, welche Aktionen möglich sind und wie sie ausgeführt werden? & 5.2.2.2 & Abweichung (teils fehlende Infotexte) & 2 \\\hline
2.2 & Lassen sich die zur Fotoerfassung und -verwaltung erforderlichen Steuerelemente zuverlässig auffinden? & 5.2.2.3 & erfüllt & 0 \\\hline
2.3 & Ist eindeutig erkennbar, welche Elemente interaktiv sind und welche nur der Anzeige dienen (z.\,B. Stammdaten als Prüfhilfe)? & 5.2.2.4 & Abweichung (unterschiedliche Button-Designs) & 1 \\\hline
2.4 & Werden Begriffe in einem den handwerklich geprägten Benutzenden vertrauten Vokabular dargestellt, insb. bei den SAP-Attributen der Meldungsauswahl? & 5.2.2.5 & Abweichung (Meldungsart in SAP-Kürzel) & 2 \\\hline
2.5 & Wird der Bearbeitungsfortschritt der Aufgabe sichtbar gemacht (z.\,B. Anzahl erfasster Fotos, Bezug zur Meldung)? & 5.2.3.1 & erfüllt & 0 \\\hline
2.6 & Werden Benutzende über Änderungen des Systemzustands informiert (z.\,B. übernommenes/verworfenes Foto, erreichte Obergrenze)? & 5.2.3.2 & erfüllt & 0 \\\hline
2.7 & Ist erkennbar, wann und wo eine Eingabe erforderlich ist (z.\,B. Pflicht zur Meldungsauswahl vor der Fotoerfassung)? & 5.2.3.3 & erfüllt & 0 \\\hline

\multicolumn{5}{|l|}{\textbf{Prinzip 3 -- Erwartungskonformität}} \\\hline
3.1 & Entsprechen die vorgesehenen Ablaufschritte dem Aufgabenverständnis der Instandhaltungsmitarbeitenden? & 5.3.2.1 & erfüllt & 0 \\\hline
3.2 & Liefert die Anwendung sofortige, angemessene Rückmeldungen auf Handlungen (z.\,B. nach Auslösen der Kamera, Übernehmen/Verwerfen)? & 5.3.2.3 & erfüllt & 0 \\\hline
3.3 & Wird bei längeren Reaktionszeiten (Laden der Meldungsliste, Foto-Upload über WLAN) eine Fortschritts- bzw. Warteinformation ausgegeben? & 5.3.2.4 & Abweichung (keine Warteinformationen) & 2 \\\hline
3.4 & Werden den Benutzenden vertraute sprachliche Konventionen (deutschsprachige Oberfläche) konsistent verwendet? & 5.3.3.1 & Abweichung (Durchmischung Deutsch/Englisch) & 3 \\\hline
3.5 & Unterscheiden sich Elemente mit unterschiedlichem Zweck (Auswahl, Aufnahme, Löschen) erkennbar voneinander? & 5.3.3.3 & erfüllt & 0 \\\hline
3.6 & Reagiert die Anwendung angemessen auf Änderungen angeschlossener Ressourcen (z.\,B. WLAN-Abbruch, nicht verfügbare Kamera)? & 5.3.4.2 & erfüllt & 0 \\\hline
3.7 & Werden Inhalte auf der genutzten Anzeigegröße (Smartphone) sowie bei Orientierungswechsel angemessen dargestellt? & 5.3.4.3 & erfüllt & 0 \\\hline

\multicolumn{5}{|l|}{\textbf{Prinzip 4 -- Erlernbarkeit}} \\\hline
4.1 & Unterstützt die Anwendung Benutzende dabei, die nötigen Funktionen selbstständig zu entdecken und zielführend zu nutzen? & 5.4.2.1 & erfüllt & 0 \\\hline
4.2 & Können Benutzende die Anwendung gefahrlos erkunden, ohne unbeabsichtigte oder nicht umkehrbare Folgen? & 5.4.3.1 & erfüllt & 0 \\\hline
4.3 & Bietet die Anwendung geeignete Alternativen für das Auffinden der richtigen Meldung (Filter sowie Freitextsuche)? & 5.4.3.2 & erfüllt & 0 \\\hline
4.4 & Liefert die Anwendung Rückmeldungen, anhand derer Benutzende die Folgen ihrer Aktionen erkennen und erlernen können? & 5.4.4.1 & erfüllt & 0\\\hline

\multicolumn{5}{|l|}{\textbf{Prinzip 5 -- Steuerbarkeit}} \\\hline
5.1 & Können Benutzende den Dokumentationsvorgang jederzeit unterbrechen? & 5.5.2.1 & erfüllt & 0 \\\hline
5.2 & Können Benutzende einen unterbrochenen Vorgang fortsetzen, ohne erfasste Fotos oder die Meldungsauswahl zu verlieren? & 5.5.2.2 & erfüllt & 0 \\\hline
5.3 & Stehen für die Bilderfassung unterschiedliche Mittel bereit (Gerätekamera oder Upload vorhandener Dateien)? & 5.5.3.2 & erfüllt & 0 \\\hline
5.4 & Lässt sich zumindest die letzte Aktion rückgängig machen (Foto verwerfen oder übernommenes Foto wieder entfernen)? & 5.5.3.4 & erfüllt & 0 \\\hline

\multicolumn{5}{|l|}{\textbf{Prinzip 6 -- Robustheit gegen Benutzungsfehler}} \\\hline
6.1 & Verzichtet die Anwendung darauf, bereits im SAP-System vorhandene Informationen manuell eingeben zu lassen? & 5.6.2.1 & erfüllt & 0 \\\hline
6.2 & Erfolgt die Meldungsauswahl über eine Suchhilfe (Wiedererkennung) statt manueller Eingabe der Meldungsnummer (Erinnerung)? & 5.6.2.2 & erfüllt & 0 \\\hline
6.3 & Lässt die Meldungsauswahl nur aktuell gültige bzw. zulässige Meldungen zu? & 5.6.2.3 & Abweichung (Filter schließen sich nicht aus) & 2 \\\hline
6.4 & Ist sichergestellt, dass erfasste Fotos nicht durch Benutzungs- oder Systemfehler (z.\,B. WLAN-Abbruch beim Upload) verloren gehen? & 5.6.2.4 & erfüllt & 0\\\hline
6.5 & Werden ungeeignete Dateien (z.\,B. unzulässiges Format beim Upload) vor der Weiterverarbeitung erkannt? & 5.6.2.5 & Abweichung (keine Fehlermeldung bei falschem Dateityp) & 3 \\\hline
6.6 & Fordert die Anwendung vor Handlungen mit nicht trivialen Auswirkungen (z.\,B. Löschen eines Fotos) eine Bestätigung an? & 5.6.2.6 & Abweichung (Foto wird ohne Rückfrage mit einem Klick gelöscht)& 2 \\\hline
6.7 & Unterstützt die Anwendung dabei, eine Fehlauswahl (z.\,B. falsche Meldung) über die angezeigten Stammdaten zu erkennen und zu korrigieren? & 5.6.4.1 & erfüllt & 0 \\\hline
6.8 & Gibt die Anwendung präzise, verständliche und höflich formulierte Fehlermeldungen aus, die ohne SAP-Fachkenntnisse verständlich sind? & 5.6.4.2 & Abweichung (Popup erscheint, enthält jedoch nur den Hinweis \enquote{Fehlermeldung kann nicht geladen werden} statt des eigentlichen Fehlertextes) & 3 \\\hline
6.9 & Liefern Fehlermeldungen konstruktive Erläuterungen, die die Behebung erleichtern? & 5.6.4.3 & Abweichung (Meldung nennt keine Handlungsanweisung zur Behebung) & 3 \\\hline
6.10 & Werden die zur Fehlerkorrektur erforderlichen Schritte und der Aufwand gering gehalten? & 5.6.4.4 & erfüllt & 0 \\\hline

\multicolumn{5}{|l|}{\textbf{Prinzip 7 -- Benutzerbindung}} \\\hline
7.1 & Bestätigt die Anwendung die erfolgreiche Erfassung und Zuordnung der Fotos und macht erkennbar, dass keine offenen Schritte verbleiben? & 5.7.2.1 & erfüllt & 0 \\\hline
7.2 & Antizipiert die Anwendung wiederkehrende Erfordernisse (z.\,B. automatische, an die Meldungsnummer gekoppelte Benennung und Zuordnung)? & 5.7.2.4 & erfüllt & 0 \\\hline
7.3 & Stellt die Anwendung auch unter erschwerten Bedingungen vor Ort (Handschuhe, gelegentliche Nutzung) keine unverhältnismäßigen Anforderungen? & 5.7.2.7 & Abweichung (kein Handschuhmodus) & 1 \\\hline
7.4 & Vermittelt die Anwendung berechtigtes Vertrauen, dass die Fotos korrekt zugeordnet und zuverlässig abgelegt werden? & 5.7.3.1/.2 & erfüllt & 0 \\\hline

\end{longtable}
}


